% ----------------------------------------------------------
% Running example schema for this section:
% Artist(id, name, country)
% Album(id, title, year, genre, artist_id)
% Track(id, title, duration_sec, album_id)
% ----------------------------------------------------------

\begin{frame}{Joins -- Overview}
  SQL \textbf{JOINs} combine rows from two or more tables based on a related column.
  They are the SQL realisation of relational-algebra join operators and are essential
  whenever data is spread across multiple tables -- which is always the case in a
  normalised schema.

  \begin{block}{Join types covered in this section}
    \begin{itemize}
      \item \textbf{INNER JOIN} -- only matching rows from both sides
      \item \textbf{NATURAL JOIN} -- automatic column-name matching
      \item \textbf{LEFT (OUTER) JOIN} -- all rows from the left table
      \item \textbf{RIGHT (OUTER) JOIN} -- all rows from the right table
      \item \textbf{Left Anti-Join} -- rows in A with no match in B
      \item \textbf{Right Anti-Join} -- rows in B with no match in A
      \item \textbf{FULL OUTER JOIN} -- all rows from both tables
    \end{itemize}
  \end{block}

  \begin{exampleblock}{Running Example -- Music Schema}
    \texttt{Artist(id, name, country)} \quad
    \texttt{Album(id, title, year, artist\_id)} \quad
    \texttt{Track(id, title, duration\_sec, album\_id)}
  \end{exampleblock}
\end{frame}

% ----------------------------------------------------------
\begin{frame}{Join Types -- Visual Overview}
  \begin{center}
  \renewcommand{\arraystretch}{1.6}
  \begin{tabular}{@{}llll@{}}
    \toprule
    \textbf{Join Type}  & \textbf{Venn} & \textbf{Returns} & \textbf{NULLs?} \\
    \midrule
    INNER JOIN       & \vennjoin{0}{1}{0} & Matching rows only           & No \\
    LEFT JOIN        & \vennjoin{1}{1}{0} & All A + matched B            & B cols \\
    RIGHT JOIN       & \vennjoin{0}{1}{1} & All B + matched A            & A cols \\
    Left Anti-Join   & \vennjoin{1}{0}{0} & A rows with no match in B    & B cols \\
    Right Anti-Join  & \vennjoin{0}{0}{1} & B rows with no match in A    & A cols \\
    FULL OUTER JOIN  & \vennjoin{1}{1}{1} & All rows from both tables    & Both \\
    \bottomrule
  \end{tabular}
  \end{center}

  \begin{block}{How to read the Venn diagrams}
    Circle \textbf{A} = left table \quad Circle \textbf{B} = right table.
    The \textcolor{blue!50}{\textbf{blue shaded region}} shows rows that appear in the result.
    Unshaded regions are discarded.
  \end{block}
\end{frame}

% ----------------------------------------------------------
\begin{frame}[fragile]{INNER JOIN}
  \begin{center}\vennjoin{0}{1}{0}\end{center}

  \begin{lstlisting}[style=sqlstyle]
-- All tracks with their album title
SELECT t.title AS track, al.title AS album
FROM   Track AS t
INNER JOIN Album AS al ON t.album_id = al.id;
  \end{lstlisting}

  \begin{block}{What it does}
    Returns \textbf{only rows where the join condition is satisfied on both sides}.
    Rows in \texttt{Track} with no matching \texttt{Album} are discarded -- and vice versa.
  \end{block}

  \begin{itemize}
    \item \texttt{JOIN} without a qualifier is identical to \texttt{INNER JOIN} in all major DBMS.
    \item \textbf{Most commonly used join} in everyday SQL work.
    \item \texttt{NULL} join columns are \emph{never} matched -- \texttt{NULL = NULL} evaluates to
          \texttt{UNKNOWN}, not \texttt{TRUE}.
  \end{itemize}
\end{frame}

% ----------------------------------------------------------
\begin{frame}[fragile]{NATURAL JOIN}
  \begin{lstlisting}[style=sqlstyle]
-- Automatically joins on all columns sharing the same name
SELECT * FROM Track NATURAL JOIN Album;
-- Implicit condition: Track.album_id = Album.album_id
-- (only works if both tables have exactly "album_id")
  \end{lstlisting}

  \begin{block}{What it does}
    Joins on \textbf{all columns sharing the same name} in both tables; no \texttt{ON} clause needed.
    Duplicate join columns appear only once in the result.
  \end{block}

  \begin{alertblock}{Danger -- Silent Breakage}
    If a developer later adds a column whose name accidentally matches a column in the other table,
    the join condition changes \emph{silently} and queries return wrong results.
    Always prefer an explicit \texttt{INNER JOIN \ldots ON} in production code.
  \end{alertblock}
\end{frame}

% ----------------------------------------------------------
\begin{frame}[fragile]{LEFT (OUTER) JOIN}
  \begin{center}\vennjoin{1}{1}{0}\end{center}

  \begin{lstlisting}[style=sqlstyle]
-- All artists, including those with no albums yet
SELECT ar.name, al.title AS album
FROM   Artist AS ar
LEFT JOIN Album AS al ON ar.id = al.artist_id;
-- Artists with no albums appear with NULL in "album"
  \end{lstlisting}

  \begin{block}{What it does}
    Returns \textbf{all rows from the left table (Artist)}, plus matching rows from Album.
    Where no match exists, Album columns are \texttt{NULL}.
  \end{block}

  \begin{itemize}
    \item \texttt{OUTER} keyword is optional: \texttt{LEFT JOIN} = \texttt{LEFT OUTER JOIN}.
    \item Result always contains at least as many rows as the left table.
    \item The most frequently used outer join in practice.
  \end{itemize}
\end{frame}

% ----------------------------------------------------------
\begin{frame}[fragile]{RIGHT (OUTER) JOIN}
  \begin{center}\vennjoin{0}{1}{1}\end{center}

  \begin{lstlisting}[style=sqlstyle]
-- All albums, including those whose artist_id is missing
SELECT ar.name, al.title AS album
FROM   Artist AS ar
RIGHT JOIN Album AS al ON ar.id = al.artist_id;
-- Albums with no Artist match appear with NULL for ar.name
  \end{lstlisting}

  \begin{block}{What it does}
    Returns \textbf{all rows from the right table (Album)}, plus matching Artist rows.
    Where no match exists, Artist columns are \texttt{NULL}.
  \end{block}

  \begin{itemize}
    \item Every \texttt{RIGHT JOIN} can be rewritten as a \texttt{LEFT JOIN} by swapping table order:
          \texttt{A RIGHT JOIN B} $\equiv$ \texttt{B LEFT JOIN A}.
    \item Most developers use only \texttt{LEFT JOIN} for consistency.
  \end{itemize}
\end{frame}

% ----------------------------------------------------------
\begin{frame}[fragile]{Left Anti-Join (NULL Trick)}
  \begin{center}\vennjoin{1}{0}{0}\end{center}

  \begin{lstlisting}[style=sqlstyle]
-- Artists who have released NO albums
SELECT ar.name
FROM   Artist AS ar
LEFT JOIN Album AS al ON ar.id = al.artist_id
WHERE  al.id IS NULL;
  \end{lstlisting}

  \begin{block}{What it does}
    Returns only rows from \texttt{Artist} that have \textbf{no matching row in Album}.
  \end{block}

  \begin{itemize}
    \item The \texttt{LEFT JOIN} produces \texttt{NULL} in Album columns for unmatched artists.
    \item \texttt{WHERE al.id IS NULL} keeps only those unmatched rows.
    \item SQL has no \texttt{ANTI JOIN} keyword -- this two-step pattern is the standard idiom.
    \item Alternative: \texttt{NOT EXISTS (\ldots)} subquery -- equivalent but sometimes slower.
  \end{itemize}
\end{frame}

% ----------------------------------------------------------
\begin{frame}[fragile]{Right Anti-Join (NULL Trick)}
  \begin{center}\vennjoin{0}{0}{1}\end{center}

  \begin{lstlisting}[style=sqlstyle]
-- Albums that have no matching artist (orphaned data)
SELECT al.title
FROM   Artist AS ar
RIGHT JOIN Album AS al ON ar.id = al.artist_id
WHERE  ar.id IS NULL;
  \end{lstlisting}

  \begin{block}{What it does}
    Returns only rows from \texttt{Album} that have \textbf{no matching row in Artist}.
  \end{block}

  \begin{itemize}
    \item Mirror image of the left anti-join.
    \item Can also be expressed as a \texttt{LEFT JOIN} with tables swapped:\\
          \texttt{Album LEFT JOIN Artist ON \ldots WHERE ar.id IS NULL}
    \item Useful for finding orphaned or inconsistent data.
  \end{itemize}
\end{frame}

% ----------------------------------------------------------
\begin{frame}[fragile]{FULL OUTER JOIN}
  \begin{center}\vennjoin{1}{1}{1}\end{center}

  \begin{lstlisting}[style=sqlstyle]
-- MySQL: simulate FULL OUTER JOIN with UNION
SELECT ar.name, al.title
FROM   Artist AS ar
LEFT JOIN Album AS al ON ar.id = al.artist_id
UNION
SELECT ar.name, al.title
FROM   Artist AS ar
RIGHT JOIN Album AS al ON ar.id = al.artist_id
WHERE  ar.id IS NULL;
  \end{lstlisting}

  \begin{block}{What it does}
    Returns \textbf{all rows from both tables}. Non-matching rows get \texttt{NULL} on the other side.
  \end{block}

  \begin{itemize}
    \item PostgreSQL and SQL Server support \texttt{FULL OUTER JOIN} directly.
    \item MySQL requires the \texttt{UNION} workaround shown above.
    \item The \texttt{WHERE ar.id IS NULL} in the second part avoids duplicating matched rows.
  \end{itemize}
\end{frame}

% ----------------------------------------------------------
\begin{frame}[fragile]{Multi-Table JOIN Example}
  \textbf{Task:} List every track, its album title, and the artist name -- only for albums
  released after 2015.

  \begin{lstlisting}[style=sqlstyle]
SELECT ar.name   AS artist,
       al.title  AS album,
       t.title   AS track,
       t.duration_sec
FROM   Track   AS t
JOIN   Album   AS al ON t.album_id   = al.id
JOIN   Artist  AS ar ON al.artist_id = ar.id
WHERE  al.year > 2015
ORDER BY ar.name, al.title, t.title;
  \end{lstlisting}

  \begin{block}{Key Points}
    \begin{itemize}
      \item Chain multiple \texttt{JOIN}s one after the other -- each \texttt{ON} clause specifies
            that particular join's condition independently.
      \item Use table aliases (\texttt{AS ar}) to keep long queries readable.
      \item \texttt{WHERE} filters apply \emph{after} all joins are evaluated.
    \end{itemize}
  \end{block}
\end{frame}

% ----------------------------------------------------------
\begin{frame}{JOIN Comparison Table}
  \renewcommand{\arraystretch}{1.25}
  \begin{tabularx}{\linewidth}{@{}lXX@{}}
    \toprule
    \textbf{Join Type}  & \textbf{Includes unmatched left rows?} & \textbf{Includes unmatched right rows?} \\
    \midrule
    INNER JOIN       & No  & No \\
    LEFT JOIN        & \textbf{Yes} (NULL right cols) & No \\
    RIGHT JOIN       & No  & \textbf{Yes} (NULL left cols) \\
    Left Anti-Join   & Only unmatched left rows & No \\
    Right Anti-Join  & No  & Only unmatched right rows \\
    FULL OUTER JOIN  & \textbf{Yes} & \textbf{Yes} \\
    CROSS JOIN       & n/a -- produces full Cartesian product & n/a \\
    \bottomrule
  \end{tabularx}

  \begin{alertblock}{NULL and JOINs}
    A column containing \texttt{NULL} on the join key can \emph{never} match any row -- not even
    another \texttt{NULL}. Use \texttt{IS NULL} / \texttt{IS NOT NULL} to handle these cases explicitly.
  \end{alertblock}
\end{frame}
